\documentclass[aspectratio=169,dvipsnames,usenames]{beamer}
\usepackage{graphicx}
\usepackage{stmaryrd}
\usepackage[T1,T2A]{fontenc}
\usepackage[utf8]{inputenc}
\usepackage[english,russian]{babel}
\usepackage{amsmath}
\usepackage{amsfonts}
\usepackage{amssymb}
\usepackage{makeidx}
\usepackage{verbatim}
\usepackage{amsthm}
%\usepackage{bnf}
\usepackage{tikz}
\usepackage{enumerate}
\usepackage{mathtext}
\usepackage{mathtools}
\usepackage{mathabx}
%\usepackage[left=2cm,right=2cm,top=2cm,bottom=2cm,bindingoffset=0cm]{geometry}
\usepackage{proof}
%\usepackage{paracol}
%\usepackage{enumitem}
\usepackage{xcolor}
\usepackage{colortbl}
%\usepackage{minted}
%\usepackage{hyperref}
\setbeamertemplate{navigation symbols}{}
\usetikzlibrary{graphs}
\usetikzlibrary{graphs.standard}
\usetikzlibrary{automata,positioning}
\usepackage{float}

\begin{document}

%\theoremstyle{dfn}
\newtheorem{dfn}{Определение}[section]
\newtheorem{nte}{Замечание}[section]

\newtheorem{axiom}{Аксиома}[section]
\newtheorem{thm}{Теорема}[section]
\newtheorem{lmm}[theorem]{Лемма}
\newtheorem{statement}{Утверждение}[section]
\newtheorem{oun_paragraph}{Пункт}[section]
\newtheorem{cons}{Следствие}[section]
\newtheorem*{exm}{Пример}

\newcommand{\comb}[1]{\operatorname{\mathcal{#1}}}
\newcommand{\func}[1]{\operatorname{#1}}
\newcommand{\reduction}[1]{{\color{OrangeRed}#1}}
\newcommand{\set}[1]{\left\{#1\right\}}

\def\from#1{\par \parbox{0.7\textwidth}{\par \hfill\raggedleft \it #1}} 

\begin{frame}{}
\begin{center}
{\LARGE Линейные и уникальные типы}
\end{center}
\end{frame}

\begin{frame}{Контекст требует формализации}

Напомним правила вывода:
	$$\infer{\Gamma, \theta \vdash \theta}{}
		\quad\infer{\Gamma \vdash \varphi \rightarrow \psi}{\Gamma, \varphi \vdash \psi}\quad
		\infer{\Gamma \vdash \psi}{\Gamma \vdash \varphi \to \psi && \Gamma \vdash \varphi}$$
	

Что такое контекст?
\begin{enumerate}
\item Это множество? Ведь $\alpha, \alpha\rightarrow\beta = \alpha\rightarrow\beta,\alpha$

	$$\infer{\alpha\rightarrow\beta, \alpha \vdash \beta}
                {\infer{\alpha, \alpha\rightarrow\beta \vdash \alpha\rightarrow\beta}{}\quad\quad\infer{ \alpha\rightarrow\beta,\alpha \vdash  \alpha}{}}$$
\item Разве это множество? Ведь $\alpha,\alpha \ne \alpha$.

	$$\infer{\vdash\alpha\rightarrow\alpha\rightarrow\alpha}{\infer{\alpha\vdash\alpha\rightarrow\alpha}{\infer{\alpha,\alpha\vdash\alpha}{}}}$$
\end{enumerate}

%Требуется формализация.
\end{frame}

\begin{frame}{Структурные правила}
Перестановка:
$$\infer{\Xi,\Sigma,\Delta,H\vdash\varphi}{\Xi,\Delta,\Sigma,H\vdash\varphi}$$
Сжатие:
$$\infer{\Xi,\delta\vdash\varphi}{\Xi,\delta,\delta\vdash\varphi}$$
Ослабление:
$$\infer{\Xi,\delta\vdash\varphi}{\Xi\vdash\varphi}$$

Сжатие и ослабление предполагают содержательные действия.

$$\infer{y:\beta\vdash\lambda x^\alpha.y : \alpha\rightarrow\beta}{\infer{y:\beta,x:\alpha\vdash y : \beta}{y:\beta\vdash y : \beta}}\quad\quad
  \infer{x:\alpha\vdash ( x,x ) : \alpha\times\alpha}{\infer{x:\alpha,x:\alpha\vdash ( x,x ) : \alpha \times \alpha}{x:\alpha\vdash x:\alpha\quad\quad x:\alpha\vdash x:\alpha}}$$
\end{frame}

\begin{frame}{Два варианта удаления конъюнкции}
Вариант 1 (классический):
$$\infer{\Gamma\vdash\alpha}{\Gamma\vdash\alpha\times\beta}\quad\infer{\Gamma\vdash\beta}{\Gamma\vdash\alpha\times\beta}$$

Вариант 2 (альтернативный):
$$\infer{\Gamma,\Delta\vdash\varphi}{\Gamma\vdash\alpha\times\beta\quad\Delta,\alpha,\beta\vdash\varphi}$$

Варианты эквивалентны при наличии структурных правил (например, классический через альтернативный):
%\infer{\Gamma\vdash \alpha \times \beta}{\infer{\Gamma,\Gamma\vdash\alpha \times \beta}{\Gamma\vdash\alpha\quad\Gamma\vdash\beta}}
%\quad\quad\quad
$$\infer{\Gamma\vdash\alpha}{\Gamma\vdash\alpha\times\beta\quad\infer{\alpha,\beta\vdash\alpha}{\infer{\alpha\vdash\alpha}{}}}
$$
\end{frame}

\begin{frame}{Изоморфизм Карри-Ховарда}
\begin{itemize}
\item Сжатие --- копирование значения
\item Ослабление --- удаление значения
\item Классическая конъюнкция:
$$\infer{\Gamma\vdash \text{fst}(P):\alpha}{\Gamma\vdash P:\alpha\times\beta}\quad\infer{\Gamma\vdash \text{snd}(P):\beta}{\Gamma\vdash P:\alpha\times\beta}$$

\item Альтернативная конъюнкция:
$$\infer{\Gamma,\Delta\vdash \text{case } P \text{ of }(x,y) \rightarrow E : \varphi}{\Gamma\vdash P:\alpha\times\beta\quad\Delta, x: \alpha, y: \beta\vdash E:\varphi}$$
\end{itemize}
\end{frame}

\begin{frame}{Линейная интуиционистская логика (вариант Филиппа Вадлера)}
Грамматика:
$$\varphi ::= X\ |\ \varphi\multimap\varphi\ |\ \varphi\otimes\varphi\ |\ \varphi\ \&\ \varphi\ |\ \varphi\oplus\varphi\ |\ |\ !\varphi$$

Связки носят специальные названия: линейная импликация, мультипликативная и аддитивная конъюнкция, аддитивная дизъюнкция, фабрика (<<точно>>, <<конечно>>).

В линейной классической логике ещё рассматривают связки $\varphi\bindnasrepma\varphi$ и $?\varphi$, в изоморфизме Карри-Ховарда они не используются.

Два типа контекстов:
$\langle \alpha \rangle$ --- линейный,
$[\beta]$ --- интуиционистский

Структурные правила:

$$\infer{\Xi,\Delta,\Gamma,H\vdash\alpha}{\Xi,\Gamma,\Delta,H\vdash\alpha}\quad\quad
  \infer{\Gamma,[\alpha]\vdash\beta}{\Gamma,[\alpha],[\alpha]\vdash\beta}\quad\quad
  \infer{\Gamma,[\alpha]\vdash\beta}{\Gamma\vdash\beta}$$

Аксиомы:

$$\infer{[\alpha]\vdash\alpha}{}\quad\quad\infer{\langle\alpha\rangle\vdash\alpha}{}$$


\end{frame}

\begin{frame}{Правила для связок}
Правила для <<конечно>> (возможно тиражировать построение $\alpha$):

$$\infer{[\Gamma]\vdash !\alpha}{[\Gamma]\vdash\alpha}\quad\quad\infer{\Gamma,\Delta\vdash\beta}{\Gamma\vdash!\alpha\quad\quad\Delta,[\alpha]\vdash\beta}$$

Линейная импликация:
$$\infer{\Gamma\vdash\alpha\multimap\beta}{\Gamma,\langle\alpha\rangle\vdash\beta}\quad\quad\infer{\Gamma,\Delta\vdash\beta}{\Gamma\vdash\alpha\multimap\beta\quad\quad\Delta\vdash\alpha}$$
\end{frame}

\begin{frame}{Правила для связок: конъюнкция и дизъюнкция}
Мультипликативная конъюнкция (возможно построить и $\alpha$ и $\beta$ одновременно):
$$\infer{\Gamma,\Delta\vdash\alpha\otimes\beta}{\Gamma\vdash\alpha\quad\quad\Delta\vdash\beta}\quad\quad\infer{\Gamma,\Delta\vdash\varphi}{\Gamma\vdash\alpha\otimes\beta\quad\quad\Delta,\langle\alpha\rangle,\langle\beta\rangle\vdash\varphi}$$

Аддитивная конъюнкция (возможно построить $\alpha$ и возможно построить $\beta$, что-то одно по нашему выбору):
$$\infer{\Gamma\vdash\alpha\&\beta}{\Gamma\vdash\alpha\quad\quad\Gamma\vdash\beta}\quad\quad\infer{\Gamma\vdash\alpha}{\Gamma\vdash\alpha\&\beta}\quad\quad\infer{\Gamma\vdash\beta}{\Gamma\vdash\alpha\&\beta}$$

Аддитивная дизъюнкция (возможно построить $\alpha$ или возможно построить $\beta$, что-то одно по их выбору):
$$\infer{\Gamma\vdash\alpha\oplus\beta}{\Gamma\vdash\alpha}\quad\quad\infer{\Gamma\vdash\alpha\oplus\beta}{\Gamma\vdash\beta}\quad\quad
  \infer{\Gamma,\Delta\vdash\varphi}{\Gamma\vdash\alpha\oplus\beta\quad\quad\Delta,\langle\alpha\rangle\vdash\varphi\quad\quad\Delta,\langle\beta\rangle\vdash\varphi}$$

\end{frame}

\begin{frame}{Пример: интуиционистская импликация}
Введём обозначение: $$\alpha\rightarrow\beta := !\alpha\multimap\beta$$

Заметим, что такая импликация ведёт себя как интуиционистская:
$$\infer{\Gamma\vdash !\alpha\multimap\beta}{\infer{\Gamma,\langle!\alpha\rangle\vdash\beta}{\langle!\infer{\alpha\rangle\vdash!\alpha}{}\quad\quad\Gamma,[\alpha]\vdash\beta}}$$

$$\infer{\Gamma,[\Delta]\vdash\beta}{\Gamma\vdash!\alpha\multimap\beta\quad\quad\infer{[\Delta]\vdash!\alpha}{[\Delta]\vdash\alpha}}$$
\end{frame}

\begin{frame}{Вложение остальных интуиционистских связок в линейные}
Например, можно так:
\begin{center}\begin{tabular}{l}
$\alpha\rightarrow\beta := !\alpha\multimap\beta$\\
$\alpha\times\beta := \alpha\&\beta$\\
$\alpha+\beta := !\alpha\oplus!\beta$
\end{tabular}\end{center}

\end{frame}

\begin{frame}{Комбинаторный базис BCKW}

\begin{itemize}
\item $B := \lambda x.\lambda y.\lambda z.x\ (y\ z)$ (комбинатор $Z$, Zusammensetzung)
\item $C := \lambda x.\lambda y.\lambda z.x\ z\ y$ (комбинатор $T$, verTauschnung)
\item $K := \lambda x.\lambda y.x$ (Konstanz)
\item $W := \lambda x.\lambda y.x\ y\ y$
\end{itemize}
$BC$ --- линейная логика (значение нельзя ни удалить, ни скопировать)

$BCK$ --- аффинная логика (значение можно удалить, но не скопировать)

$BCKW$ --- интуиционистская логика
\vspace{1cm}

Почему? Например, $W:(\alpha\rightarrow\alpha\rightarrow\beta)\rightarrow(\alpha\rightarrow\beta)$.

\end{frame}

\begin{frame}{}
\begin{center}
{\LARGE Реализация\\Уникальные типы}
\end{center}
\end{frame}

\begin{frame}{Язык}
Роды:
$$\kappa ::= \mathcal{T}\ |\ \mathcal{U}\ |\ \star\ |\ \kappa\rightarrow\kappa$$

Типы:
\begin{center}\begin{tabular}{ll}
$\texttt{Int}$ & $: \mathcal{T}$\\
$\rightarrow$ & $: \star\rightarrow\star\rightarrow\mathcal{T}$\\
$\circ,\times$ & $: \mathcal{U}$\\
$\vee,\wedge$ & $: \mathcal{U}\rightarrow\mathcal{U}\rightarrow\mathcal{U}$\\
$\neg$ & $: \mathcal{U}\rightarrow\mathcal{U}$\\
$\texttt{Attr}$ & $: \mathcal{T}\rightarrow\mathcal{U}\rightarrow\star$
\end{tabular}\end{center}

Обозначения:
\begin{center}\begin{tabular}{l}
$t^u := \texttt{Attr}\ t\ u$\\
$a \rightarrow^u b := \texttt{Attr}\ (a\rightarrow b)\ u$
\end{tabular}\end{center}

Значения:
$$E ::= x^\odot\ |\ x^\otimes\ |\ \lambda x.E\ |\ E\ E$$
\end{frame}

\begin{frame}{Типизация}
Правила типизации имеют такой вид:
$$\Gamma\vdash E:\tau|_{fv}$$

Здесь $\Gamma$ сопоставляет переменным типы рода $\star$.
$fv$ сопоставляет переменным типы рода $\mathcal{U}$.

Правила вывода:
$$\infer{\Gamma,x:t^u\vdash x^\odot:t^u|_{x:u}}{}\quad\quad\infer{\Gamma,x:t^\times\vdash x^\otimes:t^\times|_{x:\times}}{}$$
$$\infer[fv' := fv\setminus\{x:\tau\}]{\Gamma\vdash\lambda x.E:a \rightarrow^{\bigvee fv'} b|_{fv'}}{\Gamma,x:a\vdash E:b|_{fv}}\quad\quad
  \infer{\Gamma\vdash E_1\ E_2:b|_{fv_1\cup fv_2}}{\Gamma\vdash E_1:a\rightarrow^u b|_{fv_1}\quad\quad\Gamma\vdash E_2:a|_{fv2}}$$
\end{frame}

\begin{frame}[fragile]{Смысл булевский выражений}
Будем считать уникальность истиной, а неуникальность --- ложью.
И рассмотрим, например, функцию \verb!fst!:

$$\texttt{fst} : (t^u, s^v)^{w\vee u} \rightarrow t^u$$

Это то же самое, что и

$$\texttt{fst} : (t^u, s^v)^w \rightarrow t^u, w \le u$$

Однако, подобные выражения могут быть разрешены с помощью \emph{булевской унификации}.
\end{frame}

\begin{frame}[fragile]{Где применяются линейные/уникальные типы}
Ручное распределение памяти:
\begin{verbatim}
void compute() {
    char* x = new char[1024];
    char* y = x;
    y[10] = 'a';
    delete [] x;      // delete y; нельзя! будут ошибки.
}
\end{verbatim}

Автоматическое распределение памяти (сборка мусора, подсчёт ссылок):
\begin{verbatim}
public void compute() {          fn compute() {
    int[] x = new int[1024];         let x = Arc::new("abcde");
    int[] y = x;                     let y = x.clone();
    x[10] = 15;                  }
}
\end{verbatim}

А давайте посчитаем количество ссылок при компиляции. Значение линейного/аффинного типа всегда существует в единственном экземпляре.
\end{frame}

\begin{frame}[fragile]{Практический пример: Раст}
\begin{itemize}
\item В Расте (если не использовать небезопасные конструкции) ссылки
можно копировать неограниченно, ограничения в <<заимствовании>> (borrow):
\begin{enumerate}
\item На чтение (immutable, неизменяющее заимствование): количество одновременных заимствований неограничено;
\item На запись (mutable, изменяющее заимствование): только однократно.
\end{enumerate}
%Оба способа заимствования применить одновременно невозможно.
{\small
\begin{verbatim}
fn main() {
    let mut v = [1,2,3].to_vec();
    let vv = &mut v;
    if let Some(v_0) = v.get(0) {    // заимствование v для чтения
        vv.insert(2, 123);           // заимствование v для записи
        println!("v[0] = {}", v_0);
    }
}
>>> error[E0502]: cannot borrow `v` as immutable because it is 
                  also borrowed as mutable
\end{verbatim}
}
\item Значение можно <<удалить>> (drop) --- скорее похоже на аффинные типы.
\item Правила вывода задаются неформально.
\end{itemize}
\end{frame}

\begin{frame}[fragile]{Линейные и аффинные типы: отличия}
\begin{itemize}
\item С аффинными типами допустимо ослабление --- нельзя (дёшево) размножать, но можно уничтожить
(массив, вектор, хэш-таблица).

\item Внешний мир для программы (монада IO) --- значение линейного типа.
Нельзя создать, удалить, можно только изменить.
Изменение производится пользователем при вводе значения --- или компьютером, при выводе.

\item Однако, IO в Хаскеле только имитирует поведение линейного типа, поскольку монады
можно удалять, копировать и склеивать с помощью bind:
$$\begin{array}{ll}\text{return}: x \rightarrow \tau(x) & \text{(создание монады)}\\
\texttt{(>{}>{}=)}: \tau(x) \rightarrow (x \rightarrow \tau(y)) \rightarrow \tau(y) & \text{(bind, склейка монад)}\end{array}$$

\item Рассмотрим следующие смелые манипуляции с внешним миром (IO):
\small\begin{verbatim}
main :: IO ()
main = let m = return () in (m1, m2, m3) = (m, m, m) in 
       (m1 >>= \u -> putStr "Hello, ") >> (m2 >>= \u -> putStr "IO!")
\end{verbatim}
\end{itemize}
\end{frame}


\begin{frame}{}
\begin{center}
{\LARGE Объектно-ориентированное программирование и ТТ}
\end{center}
\end{frame}

\begin{frame}{Отношение <<быть подтипом>>}
\begin{dfn}
Будем говорить, что $A$ --- подтип $B$: $$A <: B$$
если значения типа $A$ --- часть множества значений типа $B$.
\end{dfn}
\begin{exm}
\begin{itemize}
\item В Java $\text{String} <: \text{Object}$.
\item В Java $\text{int} \not<: \text{Integer}$.
\end{itemize}
\end{exm}
\end{frame}

\begin{frame}{Ко- и контравариантность}
Определение взято из теории категорий (и упрощено):
\begin{dfn}
Пусть заданы два упорядоченных множества, $\langle C, < \rangle$ и $\langle D, < \rangle$.
Ковариантное (контравариантное) отображение --- $f: C \rightarrow D$, что 
\begin{itemize}
\item при $x<y$ всегда $f(x) < f(y)$ (ковариантное отображение);
\item при $x<y$ всегда $f(y) < f(x)$ (контравариантное отображение).
\end{itemize}
\end{dfn}
\end{frame}

\begin{frame}{Операции на типах}
Пусть $g : \star \rightarrow \star$. Что можно сказать о его ко- или контравариантности?

\begin{itemize}
\item Функции ковариантны по результату: $g(\sigma) := \text{int}\rightarrow\sigma$. 

Если $\alpha <: \beta$, $a : g(\alpha)$, $b : g(\beta)$ то $a(n): \alpha, b(n): \beta$, 
поэтому функцию $b$ можно всегда заменить на $a$. Отсюда,
$$\alpha <: \beta \text{ влечёт } g(\alpha) <: g(\beta)$$

\item Функции контравариантны по аргументу: $g(\sigma) := \sigma\rightarrow\text{int}$. 

Если $\alpha <: \beta$, $a : g(\alpha)$, $b : g(\beta)$, и $x : \alpha$ то
$a(x)$ и $b(x)$ законны. Отсюда,

$$\alpha <: \beta \text{ влечёт } g(\beta) <: g(\alpha)$$

\item Массивы инвариантны: $g(\alpha) := \alpha[]$.

С одной стороны, $g(\alpha).set : \alpha \rightarrow ()$, с другой стороны, $g(\alpha).get : \text{int} \rightarrow \alpha$.

В Java массивы ковариантны, что приводит к проверке времени исполнения.

\end{itemize}
\end{frame}

\begin{frame}[fragile]{Структурная и именная эквивалентность}
\begin{verbatim}
struct X { a: int, b: char } x;
struct Y { a: int, b: char } y;
\end{verbatim}

Структурная эквивалентность: $X \approx Y$

Именная эквивалентность: $X \not\approx Y$

\end{frame}

\begin{frame}{Исчисление $F_{<:}$}
Язык:
$$ T ::= x | \lambda x^\tau.T | T\ T | \lambda \alpha<:\tau.T| t\ \tau$$
Типы:
$$ \tau ::= \alpha|\top|\tau\rightarrow\tau|\forall\alpha<:\tau.\tau$$
Подтипизация:
$$\infer{\Gamma\vdash\sigma<:\sigma}{}\quad\quad\infer{\Gamma\vdash\sigma<:\tau}{\Gamma\vdash\sigma<:\rho\quad\quad\Gamma\vdash\rho<:\tau}\quad\quad\infer{\Gamma\vdash\sigma<:\top}{}$$
$$\infer{\Gamma,\alpha<:\tau\vdash\alpha<:\tau}{}\quad\quad\infer{\Gamma\vdash\sigma_1\rightarrow\sigma_2<:\tau_1\rightarrow\tau_2}{\Gamma\vdash\tau_1<:\sigma_1\quad\quad\Gamma\vdash\sigma_2<:\tau_2}$$


\end{frame}

\begin{frame}{Исчисление $F_{<:}$}
Типизация: правила системы $F$ со следующими добавлениями/изменениями
$$\infer[\text{тип. абстракция}]{\Gamma\vdash\lambda\alpha<:\tau_1.T_2 : \forall \alpha<:\tau_1.\tau_2}{\Gamma,\alpha <: \tau_1 \vdash T_2 : \tau_2}$$
$$\infer[\text{тип. применение}]{\Gamma\vdash T_1\ \tau_2 : \tau_{12}[\alpha := \tau_2]}{\Gamma\vdash T_1 : \forall\alpha <: \tau_{11}.\tau_{12}\quad\quad\Gamma \vdash \tau_2 <: \tau_{11}}$$
$$\infer[\text{приведение}]{\Gamma\vdash T:\tau}{\Gamma\vdash T:\sigma\quad\quad\Gamma\vdash\sigma<:\tau}$$
\end{frame}

\begin{frame}{Полное и ядерное правила}
Ядерное правило подтипизации:
$$\infer{\Gamma\vdash\forall\alpha<:\rho_1.\sigma_2<:\forall\alpha<:\rho_1.\tau_2}{\Gamma,\alpha<:\rho_1\vdash\sigma_2<:\tau_2}$$
Полное правило подтипизации:
$$\infer{\Gamma\vdash\forall\alpha<:\sigma_1.\sigma_2<:\forall\alpha<:\tau_1.\tau_2}{\Gamma\vdash\tau_1<:\sigma_1\quad\quad\Gamma,\alpha<:\tau_1\vdash\sigma_2<:\tau_2}$$
\end{frame}

\begin{frame}{Типизация для пар}
Можно показать:
$$\infer{\sigma_1\&\sigma_2 <: \tau_1\&\tau_2}{\Gamma\vdash\sigma_1<:\tau_1\quad\quad\Gamma\vdash\sigma_2<:\tau_2}$$
\end{frame}

\begin{frame}{Что такое класс/объект}
\begin{itemize}
\item Определим отношение подтипизации на кортежах --- это даст наследование.
$$\text{struct }x_1 : \tau_1,\dots,x_n : \tau_n\text{ end} ::= \langle x_1 : \tau_1, \langle\dots\langle x_n : \tau_n, T : \top\rangle\rangle$$
Упростим, убрав имена полей (не теряя общности):
$$\text{struct }\tau_1,\dots,\tau_n\text{ end} ::= \tau_1\&\dots(\tau_n\&\top)$$
%Можно показать, что:
%$$\infer{\sigma_1\&\sigma_2 <: \tau_1\&\tau_2}{\Gamma\vdash\sigma_1<:\tau_1\quad\quad\Gamma\vdash\sigma_2<:\tau_2}$$
Тогда $$\tau_n\&\tau_{n+1} <: \tau_n\&\top$$
Отсюда, $$\text{struct }x_1 : \tau_1,\dots,x_n : \tau_n\text{ end} <: \text{struct }x_1 : \tau_1,\dots,x_{n-1} : \tau_{n-1}\text{ end}$$
\item Приведения типов и т.п. получаются согласно общим правилам $F_{<:}$.
\item По необходимости добавим экзистенциальные типы.
\item При необходимости сделать именную эквивалентность из структурной можно, включив какие-нибудь токены в тип .
\end{itemize}
\end{frame}

\begin{frame}{Objective Caml}
\begin{itemize}
\item ML --- Meta Language, Робин Милнер, 1970е.
\item Standard ML, 1983.
\item Category Abstract Machine Language (Caml), 1985.
\item ZINC project (ZINC Is Not CAML): ``Toplevels considered harmful'', Ксавье Леруа (Xavier Leroy), 1990.
\item Objective Caml, 1996.
\end{itemize}
\end{frame}

\begin{frame}{Объекты, классы, подтипы}
\begin{itemize}
\item Объектно-ориентированность --- набор различных конструкций, которые можно собирать по-разному.
\item Объектно-ориентированность в Окамле собрана иначе, чем в Джаве или в Смолтоке.
\item В Окамле есть две различных конструкции: модули (соответствуют абстрактным типам данных), 
и объекты и классы (предназначены для построения иерархии наследования).
\item Отношение подтипирования определено независимо от модулей и классов, в том числе и на обычных типах.
\end{itemize}
\end{frame}

\begin{frame}[fragile]{Полиморфный вариантный тип}
Традиционный алгебраический тип не допускает пересечение вариантов:
\begin{verbatim}
type abc = A | B | C;;
type ac = A | C;;             (* ac.C <> abc.C *)
\end{verbatim}

Однако, есть полиморфный вариантный тип:
\begin{verbatim}
type abc = [`A | `B | `C];;
type ac = [`A | `C ];;
\end{verbatim}

Заметим, что $[>`A|`C] :> [>`A|`B|`C]$.

\vspace{0.5cm}
Автоматический вывод типов не может построить правильное подтипирование, но можно указать его явно:

\begin{verbatim}
let f a = `A;;                  (* f : 'a -> [> `A ] *)
let g = (f () : ac);;           (* ошибка: у типа f() нет тэга `C *)
let g = (f () : [`A] :> ac);;   (* g : ac, явное приведение типа *)
\end{verbatim}
\end{frame}

\begin{frame}[fragile]{Модули: ковариантный интерфейс}
\footnotesize
\begin{verbatim}
module type Writer = sig
    type +'b t
    val empty : unit -> 'b t
    val push : 'b t -> 'b -> 'b t
    val print_len : 'b t -> unit
end

module Writer : Writer = struct
    type 'b t = 'b list
    let empty () = []
    let push l x = x :: l
    let print_len l = print_int (List.length l)
end

let w = Writer.empty ();;
let w = Writer.push w `A;;
let w = Writer.push w `C;;
let u = (w : [`A|`C] Writer.t :> [`A|`C|`E] Writer.t);;
let u = Writer.push u `E;;
\end{verbatim}
%Writer.print_len u;;
\end{frame}

\begin{frame}[fragile]{Модули: контравариантный интерфейс}
\footnotesize
\begin{verbatim}
module type Counter = sig
    type -'b t
    val empty : ('b->int->int) -> 'b t
    val push : 'b t -> 'b -> 'b t
    val print_len : 'b t -> unit
end

module Counter : Counter = struct
    type 'b t = int * ('b->int->int)
    let empty f = (0,f)
    let push (l,f) x = (f x l, f)
    let print_len (l,f) = print_int l
end

let w = Counter.empty (fun x l -> l + match x with `A -> 0 | `B -> 1 | `C -> 2);;
let w = Counter.push w `A;;
let w = Counter.push w `C;;
let u = (w : [`A|`C] Counter.t :> [`A] Counter.t);; (* однако, [`A] :> [`A|`C] *)
let u = Counter.push u `A;;
\end{verbatim}
%Counter.print_len u;;
\end{frame}

\begin{frame}[fragile]{Объекты}
Зададим объект:

\footnotesize
\begin{verbatim}
type square = < area : float; width : int >;;

let square w = object
  method area = Float.of_int (w * w)
  method width = w
end;;
\end{verbatim}

\normalsize
И нечто общее:

\footnotesize
\begin{verbatim}
let coin = object
  method shape = circle 5
  method color = "silver"
end;;

let map = object
  method shape = square 10
end;;

type item = < shape : shape >;;
let items = [ (coin :> item) ; (map :> item) ];;
\end{verbatim}
\end{frame}

\begin{frame}{Формализация: подтипы}
\begin{itemize}
\item ``Быть подтипом'' определяется рекурсивно --- согласно структуре.
\item Приведение типа:
$$\infer{\Gamma\vdash(\alpha: \tau<:\tau'):\theta(\tau')}{\tau <: \tau'\quad\quad\Gamma\vdash\alpha:\theta(\tau)}$$
\end{itemize}
\end{frame}

\begin{frame}[fragile]{Формализация: экзистенциальные типы и модули}
\begin{itemize}
\item 
\begin{verbatim}
module type Counter = sig
    type -'b t
    val empty : ('b->int->int) -> 'b t
    ...
end
\end{verbatim}
То есть, $\exists \beta.E : (\beta\rightarrow N \rightarrow N) \rightarrow \tau(\beta)$
\item Заменим экзистенциальный тип (интерфейс модуля) на его каноническое представление:
$(B\rightarrow N \rightarrow N)\rightarrow \tau(B)$
\item Правила подтипирования для типа, на основании его канонического представления:
$$ \infer{\Gamma\vdash\exists\alpha.R <: \exists\alpha'.R}{\Gamma,\alpha\vdash R<: R'[\alpha':=\tau]}$$
\end{itemize}
\end{frame}

\end{document}
